\chapter[Adversarial Machine Learning]{Adversarial Machine Learning: Attacks and Defenses}

In 2014 Goodfellow, Shlens, and Szegedy showed one picture of a panda. The GoogLeNet
network called it a panda with confidence $57.7$ percent. They added a perturbation
that the eye cannot see, of size $0.007$ on the scale $[0,1]$, which is about two
units out of 255 gray levels. The same network then called that picture a gibbon with
confidence $99.3$ percent. Nothing changed for a human. What changed was only the
direction of the perturbation: every pixel was shifted a little, and the direction of
that shift was chosen from the sign of the derivative of the error with respect to the
pixel.

Numbers like that are not an oddity of one model. In this chapter we compute for
ourselves a small case that can be rechecked. A two-layer network with 64 inputs
reaches accuracy $97.99$ percent on 597 test images of digits. With a perturbation
budget $0.15$ per pixel, an iterative attack of 40 steps lowers that accuracy to
$10.72$ percent; at the same budget, a random perturbation without a direction only
lowers it to $95.48$ percent. So the direction that the gradient chooses is not
decoration: it buys the attacker about eighty points of accuracy.

Chapter 25 taught us to compute the derivative of the error with respect to the
weights, and Chapter 27 taught us to measure accuracy honestly. Neither is enough.
Accuracy is an average over one data distribution; security is a question about the
worst point inside one allowed set. A high average and a bad worst point can live
together, and that is what makes an ordinary evaluation report misleading. This
chapter builds the vocabulary for that worst-case question, computes its first attack,
then separates the defenses that raise the standard from those that only hide the
gradient.

\begin{untukanak}
Imagine a building guard who recognizes people by small habits: how many buttons of
the shirt are fastened, the bag in the left hand or the right, and hat or no hat. So
far that method has worked. Everyone entitled to enter really does differ in those
three things. Then someone comes who knows exactly how the guard sees. He undoes one
button, moves the bag to the other hand, and puts on a hat. His face is the same, his
name is the same, he is the same person. The guard still lets him in, because
according to his notes a person with those three marks is the cleaning staff.

The question that makes this story important is not ``how do we fool the guard'',
but ``how do we know that the guard can be fooled''. If we test only with random
people in random clothes, the guard will pass almost always. The error appears only
when someone deliberately chooses those three marks. Testing with chance is not the
same as testing with malice.

A computer that recognizes images works like that guard, only with thousands of small
features at once. In this chapter we will find out why a feature that looks trivial
can be chosen deliberately, why a change that the eye cannot see can flip the answer,
and how to tell a defense that is truly strong from a defense that only leaves the
attacker confused.
\end{untukanak}

\section{Threat model: what the attacker may do}
\tingkat{basic}

The sentence ``the attacker changes the input a little'' cannot be tested before we
agree on the meaning of ``a little''. Every robustness claim is always relative to
that choice, so that choice is part of the claim, not a footnote.

\kamus{threat model}{model ancaman}{The complete list of what the attacker may do:
how large the perturbation is, what it knows, and what its goal is.}
\index{threat model}

\begin{definisi}[Adversarial example and threat model]
Let $f:\mathbb{R}^{d}\to\{1,\dots,K\}$ be a classifier and
$(\mathbf{x},y)$ a correctly classified example. An
\emph{adversarial example} is $\mathbf{x}'=\mathbf{x}+\boldsymbol{\delta}$ with
$f(\mathbf{x}')\neq y$ but $\boldsymbol{\delta}$ small according to an agreed
measure. The threat model is determined by four things:
\begin{enumerate}
  \item the set of allowed perturbations
        $S(\mathbf{x})=\{\mathbf{x}+\boldsymbol{\delta}:
        \lVert\boldsymbol{\delta}\rVert\le\varepsilon\}$, where
        $\lVert\cdot\rVert$ is usually the $\ell_\infty$ norm and $\varepsilon>0$
        is the \emph{budget};
  \item the knowledge of the attacker: full access to $f$ and its gradient
        (\emph{white box}) or only to the output of $f$ (\emph{black box});
  \item the goal: any misclassification (\emph{untargeted}) or a particular
        class (\emph{targeted});
  \item other constraints, for example the perturbation must still be a valid
        object.
\end{enumerate}
The attack succeeds when there is a $\boldsymbol{\delta}\in S(\mathbf{x})$ that
changes the label.
\end{definisi}

The $\ell_\infty$ norm is chosen because it bounds every coordinate separately:
$\lVert\boldsymbol{\delta}\rVert_\infty\le\varepsilon$ means every pixel shifts by
at most $\varepsilon$. For an 8-bit image with values $0$ to $255$, the budget
$\varepsilon=1/255$ means every pixel may change by at most one level.
The $\ell_2$ norm bounds the total length of the vector and allows one pixel to
change a lot as long as the others barely change. The two answer different
questions, and a robustness claim that does not name its norm cannot be compared
with another claim.
\index{infinity norm}
\index{perturbation budget}

The most honest measure of robustness is the distance to the nearest
misclassification point,
\[
  r(\mathbf{x})=\inf\{\lVert\boldsymbol{\delta}\rVert:
  f(\mathbf{x}+\boldsymbol{\delta})\neq y\}.
\]
This quantity does not depend on the algorithm of any attacker. Unfortunately $r$
is hard to compute: finding it means maximizing a nonconvex function on a large
set. For a network with ReLU units, $r$ can be computed exactly by turning the
problem into a mixed-integer linear program, and its cost grows quickly with the
number of units. That is why almost every report in this field names what a
particular attacker managed to achieve, not the value of $r$. The gap between the
two is the source of most conclusions that are too optimistic.

One fact deserves to be recorded from the start. On tabular data, the
$\ell_\infty$ ball contains impossible points: negative ages, non-integer packet
counts, combinations of TCP flags that never appear in a real network. Examples
like that fool the model but cannot be realized, so their success tells us nothing
about the real risk. A budget that is too loose produces attacks that are
spectacular and meaningless.

\section{Why the sign of the gradient is enough}
\tingkat{intermediate}

Now we attack. Let the error of the model be expressed by a function
$L(\mathbf{x},y)$, for example cross-entropy. We want to raise that error as much
as possible while keeping every coordinate inside the box. For a linear function,
the answer can be written in closed form.

\kamus{fast gradient sign method}{metode tanda gradien}{A one-step attack that
shifts every coordinate by the budget, with the direction following the sign of the
derivative of the error.}
\index{FGSM}

\begin{lemma}[Maximum of a linear form on the $\ell_\infty$ ball]
For $\mathbf{w}\in\mathbb{R}^{d}$ and $\varepsilon>0$,
\[
  \max_{\lVert\boldsymbol{\delta}\rVert_\infty\le\varepsilon}
  \mathbf{w}^{\top}\boldsymbol{\delta}
  =\varepsilon\lVert\mathbf{w}\rVert_1,
\]
and the maximum is attained at $\delta_j=\varepsilon\,\mathrm{sign}(w_j)$ for every
coordinate with $w_j\neq0$.
\end{lemma}

\begin{proof}
Write $\mathbf{w}^{\top}\boldsymbol{\delta}=\sum_j w_j\delta_j$. Because
$|\delta_j|\le\varepsilon$,
\[
  \sum_j w_j\delta_j\le\sum_j|w_j|\,|\delta_j|\le\varepsilon\sum_j|w_j|.
\]
If $\delta_j=\varepsilon\,\mathrm{sign}(w_j)$, then $w_j\delta_j=\varepsilon|w_j|$
for every $j$, so that bound is attained. Coordinates with $w_j=0$ do not
affect the value, so their choice is free.
\end{proof}

For a smooth function, we use the first-order approximation of $L$ around
$\mathbf{x}$ together with the lemma above. The result is the recipe that
Goodfellow and colleagues call the \emph{fast gradient sign method}:
\[
  \mathbf{x}^{\mathrm{adv}}
  =\mathbf{x}+\varepsilon\,\mathrm{sign}\bigl(\nabla_{\mathbf{x}}L(\mathbf{x},y)\bigr).
\]
In everyday language it goes like this. The gradient tells us in which direction
the error rises fastest, but we may only walk a distance $\varepsilon$; only the
sign of each component is used, while the size is fixed by the budget. That is why
this method is called the sign attack. For a linear function, that step is exactly
optimal; Figure~\ref{fig:batas} shows the result on a model with two features.

\begin{figure}[htbp]
\centering
{\scriptsize
\begin{tikzpicture}
\node[at={(0.000cm,0.000cm)},inner sep=0pt,minimum size=0pt] (o) {};
\node[at={(8.400cm,0.000cm)},inner sep=0pt,minimum size=0pt] (ex) {};
\node[at={(0.000cm,8.400cm)},inner sep=0pt,minimum size=0pt] (ey) {};
\draw[-latex,black!75] (o) -- (ex);
\draw[-latex,black!75] (o) -- (ey);
\node[anchor=north,at={(4.200cm,-1.120cm)}]{upper-half ink};
\node[rotate=90,anchor=south,at={(-1.15cm,5.10cm)}]{lower-half ink};
\node[anchor=north,at={(2.380cm,-0.280cm)}]{0.20};
\node[anchor=north,at={(6.580cm,-0.280cm)}]{0.35};
\node[anchor=east,at={(-0.224cm,2.660cm)}]{0.35};
\node[anchor=east,at={(-0.224cm,6.860cm)}]{0.50};
\node[at={(1.439cm,0.000cm)},inner sep=0pt,minimum size=0pt] (b0) {};
\node[at={(8.400cm,7.224cm)},inner sep=0pt,minimum size=0pt] (b1) {};
\draw[black,line width=1.2pt] (b0) -- (b1);
\node[circle,fill=black!55,inner sep=1.35pt,at={(2.358cm,4.399cm)}]{};
\node[circle,fill=black!55,inner sep=1.35pt,at={(3.014cm,5.602cm)}]{};
\node[circle,fill=black!55,inner sep=1.35pt,at={(2.796cm,2.376cm)}]{};
\node[circle,fill=black!55,inner sep=1.35pt,at={(3.178cm,5.274cm)}]{};
\node[circle,fill=black!55,inner sep=1.35pt,at={(0.444cm,3.524cm)}]{};
\node[circle,fill=black!55,inner sep=1.35pt,at={(1.975cm,4.290cm)}]{};
\node[circle,draw=black!55,line width=0.5pt,inner sep=1.15pt,at={(5.585cm,2.649cm)}]{};
\node[circle,draw=black!55,line width=0.5pt,inner sep=1.15pt,at={(5.749cm,2.868cm)}]{};
\node[circle,draw=black!55,line width=0.5pt,inner sep=1.15pt,at={(5.147cm,1.391cm)}]{};
\node[circle,draw=black!55,line width=0.5pt,inner sep=1.15pt,at={(3.780cm,0.844cm)}]{};
\node[circle,draw=black!55,line width=0.5pt,inner sep=1.15pt,at={(6.678cm,3.032cm)}]{};
\node[circle,draw=black!55,line width=0.5pt,inner sep=1.15pt,at={(5.803cm,2.868cm)}]{};
\node[circle,draw=black,line width=1pt,inner sep=2.8pt,at={(2.303cm,7.079cm)}]{};
\node[circle,fill=black,inner sep=0.9pt,at={(2.303cm,7.079cm)}]{};
\node[rectangle,fill=accent,inner sep=2.3pt,at={(5.663cm,3.719cm)}]{};
\node[at={(2.303cm,7.079cm)},inner sep=0pt,minimum size=0pt] (x0) {};
\node[at={(5.663cm,3.719cm)},inner sep=0pt,minimum size=0pt] (xa) {};
\draw[-latex,accent,line width=1.1pt] (x0) -- (xa);
\node[anchor=south east,at={(2.135cm,7.443cm)}]{clean};
\node[anchor=north west,at={(5.943cm,3.551cm)}]{adversarial};
\node[anchor=east,at={(2.660cm,5.600cm)}]{class 6};
\node[anchor=south west,at={(5.684cm,1.260cm)}]{class 0};
\node[anchor=south east,at={(5.475cm,4.647cm)}]{boundary};
\end{tikzpicture}}
\caption{Decision boundary of a model with two features on the digits $0$ against
$6$ ($236$ training data, $123$ test, accuracy $98.37$ percent). Thick circle: a
clean example $(0.197;0.508)$ with score $+1.936$. Red square: the result of FGSM
with $\varepsilon=0.12$ at $(0.317;0.388)$ with score $-0.208$. Its arrow is the
perturbation $(0.12;-0.12)$, $\lVert\boldsymbol{\delta}\rVert_\infty=0.12$
and $\lVert\boldsymbol{\delta}\rVert_2=0.1697$.}
\label{fig:batas}
\end{figure}

\begin{teorema}[Linearization and the remainder]
Let $L:\mathbb{R}^{d}\to\mathbb{R}$ be differentiable and let its gradient be
$M$-Lipschitz on the segment between $\mathbf{x}$ and $\mathbf{x}+\boldsymbol{\delta}$.
For every $\lVert\boldsymbol{\delta}\rVert_\infty\le\varepsilon$,
\[
  \left|L(\mathbf{x}+\boldsymbol{\delta})-L(\mathbf{x})
  -\nabla L(\mathbf{x})^{\top}\boldsymbol{\delta}\right|
  \le\frac{M}{2}\lVert\boldsymbol{\delta}\rVert_2^2
  \le\frac{Md}{2}\varepsilon^2 .
\]
\end{teorema}

\begin{proof}
With the integral form of the mean value theorem,
\[
  L(\mathbf{x}+\boldsymbol{\delta})-L(\mathbf{x})
  =\int_0^1\nabla L(\mathbf{x}+t\boldsymbol{\delta})^{\top}\boldsymbol{\delta}\,dt .
\]
Subtract $\nabla L(\mathbf{x})^{\top}\boldsymbol{\delta}
=\int_0^1\nabla L(\mathbf{x})^{\top}\boldsymbol{\delta}\,dt$, then use the
Cauchy--Schwarz inequality and the Lipschitz property:
\[
  \left|\int_0^1\bigl(\nabla L(\mathbf{x}+t\boldsymbol{\delta})
  -\nabla L(\mathbf{x})\bigr)^{\top}\boldsymbol{\delta}\,dt\right|
  \le\int_0^1 Mt\lVert\boldsymbol{\delta}\rVert_2^2\,dt
  =\frac{M}{2}\lVert\boldsymbol{\delta}\rVert_2^2 .
\]
Because $\lVert\boldsymbol{\delta}\rVert_2^2\le d\,\varepsilon^2$, the second
inequality follows.
\end{proof}

For a linear model, that remainder is zero and FGSM is not an approximation but the
exact answer. Goodfellow and colleagues use the case of logistic regression to show
this. There the sign of the gradient of the error with respect to the input is
$-\mathrm{sign}(\mathbf{w})$ times the label, so that
$\mathbf{w}^{\top}\mathrm{sign}(\mathbf{w})=\lVert\mathbf{w}\rVert_1$ and the margin
of the model shrinks by exactly $\varepsilon\lVert\mathbf{w}\rVert_1$. That is why
adversarial training on logistic regression can be written as a penalty
$\zeta\bigl(y(\varepsilon\lVert\mathbf{w}\rVert_1-\mathbf{w}^{\top}\mathbf{x}-b)\bigr)$,
that is a penalty resembling an $\ell_1$ regularizer that enters the margin, not the
cost; this exact form is used in the worked example later.

What needs to be stated honestly: the theorem above justifies the \emph{direction} of
the step, not the \emph{distance} of the step. The remainder $Md\varepsilon^{2}/2$
depends on the Lipschitz constant of the gradient, and for a ReLU network that
constant can be very large. In our experiment, the predicted rise in the error from
the linear term at $\varepsilon=0.15$ is $0.0008$ nats, whereas the rise that
actually occurs is $4.862$ nats. So linearization does not explain the size of the
damage; it only explains why the chosen direction makes sense.

The other side of the story is stronger instead. A perturbation of the same size but
with random signs only lowers accuracy from $97.99$ to $95.48$ percent, whereas FGSM
lowers it to $14.91$ percent. So what is valuable is the pattern of signs, even
though the error at the starting point is almost flat, so that the length of the
gradient and its linear forecast are small.
\index{input gradient}

So the sign of the gradient is enough not because its linear term is large, but
because the sign is more resistant to saturation of the error. The number $0.0008$
against $4.862$ is a warning that linearity holds only as an argument about
direction.

\begin{bacakeras}
Two voices, read alternately with a friend. Tap the table every time the word ``a
little'' is said.

``Shift it a little.'' ``How little?'' ``A fingertip. No more.''
``Does the answer change?'' ``It changes completely. Not doubtful, not off target: changed.''
``What do our eyes see?'' ``Nothing. Exactly the same picture.''
``And if it is shifted that far in another direction?'' ``The answer becomes correct again.''
``So what is doing the work?'' ``The direction. Not the size.''
``Where does it learn the direction?'' ``It reads the slope, asks every pixel: up or
down? Then all the answers are used at once.''
``And if it misreads the slope?'' ``Nothing happens. The picture stays, the answer stays.''

Repeat three times, faster and faster: what works is the direction, not the size.
\end{bacakeras}

\section{Iterative attacks and the min-max game}
\tingkat{intermediate}

One step is a compromise. If the gradient may be computed many times, we can climb
further: a small step in the direction of the sign of the gradient, then pull back
into the box.

\begin{definisi}[Iterative attack]
Let $\Pi$ be the projection onto $S(\mathbf{x})$. The \emph{projected gradient
descent} attack (PGD) computes
\[
  \mathbf{x}_{t+1}=\Pi\Bigl(\mathbf{x}_t+\alpha\,
  \mathrm{sign}\bigl(\nabla_{\mathbf{x}}L(\mathbf{x}_t,y)\bigr)\Bigr),
  \qquad \mathbf{x}_0=\mathbf{x},
\]
with $\alpha>0$ the step size. If $\mathbf{x}_0$ is taken at random inside
$S(\mathbf{x})$, the attack is called PGD with a \emph{random start}.
\end{definisi}

\begin{resep}
Four steps to run a bounded iterative attack:
\begin{enumerate}
  \item Fix the norm, $\varepsilon$, $\alpha$, and the number of steps; record the
        clean accuracy before a single perturbation is made.
  \item Start from $\mathbf{x}$ or from a random point inside the ball, compute the
        gradient of the error with respect to the input, then take only its sign.
  \item Step: $\mathbf{x}_{t+1}=\mathbf{x}_t+\alpha\,\mathrm{sign}(\mathbf{g})$,
        project onto the $\varepsilon$ ball, then onto $[0,1]$.
  \item Repeat, then report the accuracy at several numbers of steps and several
        $\varepsilon$; compare with FGSM, which uses $\alpha=\varepsilon$.
\end{enumerate}
Warning about order: the projection is always with respect to the original
$\mathbf{x}$, not to the previous step point, and always after the gradient step.
If the reference is swapped, the ball moves along and the perturbation exceeds the
reported budget.
\end{resep}

The projection onto the $\ell_\infty$ box cut to $[0,1]^d$ has a very
simple form: clip every coordinate, one by one.

\begin{proposisi}[Projection onto the box]
For $S=\{\mathbf{z}:\lVert\mathbf{z}-\mathbf{x}\rVert_\infty\le\varepsilon\}
\cap[0,1]^{d}$,
\[
  \Pi(\mathbf{z})_j=\min\bigl\{\max\{z_j,\,x_j-\varepsilon,\,0\},\,
  x_j+\varepsilon,\,1\bigr\}
\]
for every $j$.
\end{proposisi}

\begin{proof}
The objective function $\lVert\mathbf{y}-\mathbf{z}\rVert_2^2$ separates per
coordinate, and the allowed set also separates: every $y_j$ must lie in the interval
$[\max\{0,x_j-\varepsilon\},\min\{1,x_j+\varepsilon\}]$. For one coordinate,
minimizing $(y_j-z_j)^2$ on a closed interval means taking the nearest point, that
is $z_j$ if it is inside the interval, and the nearest end of the interval if it is
outside. That is exactly the clipping formula above.
\end{proof}

With a projection of this kind, PGD adds only two things compared with FGSM:
repeated small steps and the pull back into the box. The increase in cost is
proportional to the number of steps, because every step requires one backward pass.
Figure~\ref{fig:langkah} shows three such steps on the same model.

\begin{figure}[htbp]
\centering
{\scriptsize
\begin{tikzpicture}
\node[at={(0.000cm,0.000cm)},inner sep=0pt,minimum size=0pt] (p) {};
\node[at={(8.400cm,0.000cm)},inner sep=0pt,minimum size=0pt] (px) {};
\node[at={(0.000cm,8.400cm)},inner sep=0pt,minimum size=0pt] (py) {};
\draw[-latex,black!75] (p) -- (px);
\draw[-latex,black!75] (p) -- (py);
\node[anchor=north,at={(4.200cm,-1.344cm)}]{upper-half ink};
\node[rotate=90,anchor=south,at={(-1.05cm,4.20cm)}]{lower-half ink};
\node[anchor=north,at={(1.848cm,-0.336cm)}]{0.20};
\node[anchor=north,at={(6.888cm,-0.336cm)}]{0.35};
\node[at={(2.661cm,0.000cm)},inner sep=0pt,minimum size=0pt] (c0) {};
\node[at={(8.400cm,5.956cm)},inner sep=0pt,minimum size=0pt] (c1) {};
\draw[black,line width=1.2pt] (c0) -- (c1);
\node[at={(1.756cm,6.479cm)},inner sep=0pt,minimum size=0pt] (q0) {};
\node[circle,fill=white,draw=accent,line width=0.9pt,inner sep=2.0pt,at={(1.756cm,6.479cm)}]{};
\node[at={(3.100cm,5.135cm)},inner sep=0pt,minimum size=0pt] (q1) {};
\node[circle,fill=white,draw=accent,line width=0.9pt,inner sep=2.0pt,at={(3.100cm,5.135cm)}]{};
\node[at={(4.444cm,3.791cm)},inner sep=0pt,minimum size=0pt] (q2) {};
\node[circle,fill=white,draw=accent,line width=0.9pt,inner sep=2.0pt,at={(4.444cm,3.791cm)}]{};
\node[at={(5.788cm,2.447cm)},inner sep=0pt,minimum size=0pt] (q3) {};
\node[circle,fill=white,draw=accent,line width=0.9pt,inner sep=2.0pt,at={(5.788cm,2.447cm)}]{};
\draw[-latex,accent,line width=1.0pt,shorten >=5pt,shorten <=5pt] (q0) -- (q1);
\draw[-latex,accent,line width=1.0pt,shorten >=5pt,shorten <=5pt] (q1) -- (q2);
\draw[-latex,accent,line width=1.0pt,shorten >=5pt,shorten <=5pt] (q2) -- (q3);
\node[anchor=south west,at={(1.891cm,6.613cm)}]{$\mathbf{x}_0$};
\node[anchor=south west,at={(3.235cm,5.269cm)}]{$\mathbf{x}_1$};
\node[anchor=south west,at={(4.579cm,3.925cm)}]{$\mathbf{x}_2$};
\node[anchor=south west,at={(5.923cm,2.581cm)}]{$\mathbf{x}_3$};
\node[anchor=south east,at={(5.616cm,3.444cm)}]{boundary};
\end{tikzpicture}}
\caption{Three sign-of-gradient steps from $\mathbf{x}_0$, each step
$\alpha=0.04$. The $\lVert\boldsymbol{\delta}\rVert_\infty$ norm: $0.04$,
$0.08$, $0.12$; the $\lVert\boldsymbol{\delta}\rVert_2$ norm: $0.0566$,
$0.1131$, $0.1697$. The score of class $6$: $+1.936$, $+1.221$, $+0.506$,
$-0.208$; the boundary is crossed at the third step.}
\label{fig:langkah}
\end{figure}

There is one warning that is often missed. FGSM is PGD with one step and
$\alpha=\varepsilon$, not PGD with one step and $\alpha=\varepsilon/10$.
In our experiment, one step with $\alpha=\varepsilon/10$ only lowers
accuracy to $94.64$ percent, whereas FGSM lowers it to $14.91$
percent. So ``one step'' is not the name of one algorithm; it only means something
once its step size is named. A fair comparison uses the same step size; the
following table is computed at $\varepsilon=0.15$ with a random start.

\begin{center}
\begin{tabular}{rr}
\hline
steps & accuracy (\%) \\
\hline
1  & $94.64$ \\
2  & $88.78$ \\
5  & $66.67$ \\
10 & $24.79$ \\
20 & $11.22$ \\
40 & $10.72$ \\
\hline
\end{tabular}
\end{center}

That curve falls sharply up to ten steps, then is almost flat; adding steps after
that point rarely changes the conclusion, and that is why PGD is used as the default
attacker. Madry and colleagues read it as a sign that the inner problem is easy
enough to solve with first-order methods, even though in general it is nonconcave.

\begin{tangan}
Materials: ten cards, a marker, and a ruler. On each card write one number between
zero and one, two digits after the decimal point. Make two cards whose numbers differ
by only 0.01, for example 0.62 and 0.63.

Extra rule: when two numbers differ by less than 0.02, the higher card in the pile
wins. Place the two nearly tied cards in the top two positions, then call the topmost
card the answer of the model.

Now shift the card 0.62 up by one millimeter, until it passes the card 0.63. No
number changes, no card is swapped. Look at the topmost card: the answer changes,
because the tie-breaking rule now chooses the other card.

What must be visible: the answer changes without a single number changing, because
only the difference decides, and a thin difference can be flipped by a shift that the
eye cannot see.
\end{tangan}

From the point of view of training, attack and defense are two sides of one problem.
If the attacker maximizes the error inside the box, then the careful trainer
minimizes that worst error:
\[
  \min_{\boldsymbol{\theta}}\;
  \mathbb{E}_{(\mathbf{x},y)\sim\mathcal{D}}
  \Bigl[\max_{\boldsymbol{\delta}\in S(\mathbf{x})}
  L(\boldsymbol{\theta},\mathbf{x}+\boldsymbol{\delta},y)\Bigr].
\]
This formulation is a saddle problem: a maximum inside, a minimum outside. The inner
term is the problem of the attacker, the outer term is the problem of ordinary
training. This is \emph{adversarial training}. If the inner maximum is approximated
with FGSM, we get cheap training; if it is approximated with PGD of several steps, we
get training that is more expensive and usually more robust.
\index{adversarial training}\index{PGD}

What that formulation can and cannot promise must be separated. If the saddle value
is really attained and its value is zero, no perturbation inside the box can fool the
model, so we have a guarantee. But that saddle problem is not concave inside and not
convex outside, so any algorithm only gives an approximation: a small saddle value is
not necessarily attained, and the value that is obtained is not necessarily the
saddle value. That is why adversarial training produces a model that is \emph{more
resistant to the attacker used during training}, not a model that is proved resistant
to all attackers.

Our numbers show both sides. Model $D$, trained with FGSM at
$\varepsilon=0.15$, loses $1.51$ points of clean accuracy: from $97.99$ to
$96.48$ percent. In exchange, at $\varepsilon=0.15$ its accuracy under
FGSM is $66.67$ percent and under PGD of 40 steps $62.65$ percent, whereas
the ordinary model is left with $14.91$ and $10.72$ percent. The trade is real, but
$62.65$ percent is also not security; it only means the attacker must work
harder.

\section{Transfer, black box, and targeted attacks}
\tingkat{advanced}

Up to here our attacker knows everything. In the real world it can often only send
an input and look at the answer. What saves us when there is no gradient is not the
secrecy of the model, but the fact that the direction of attack tends to carry over
between models.

\kamus{transferability}{transferabilitas}{The property of an adversarial example that
is built for one model but also fools another model that was trained separately.}
\index{transferability}

This phenomenon was reported by Szegedy and colleagues in 2014: an example that fools
one network often fools another network with a different architecture and different
training data. Goodfellow and colleagues record its quantitative version on MNIST: an
example built from an old model gives an error of $19.6$ percent on a model trained
adversarially, whereas an example from the new model gives an error of $40.9$ percent
on the old model. The same direction, a different strength.

In our experiment, model $A$ (64 hidden units) and model $B$ (128 units) are trained
on the same training data but with different initialization. At
$\varepsilon=0.15$, PGD of 40 steps built on $A$ lowers the accuracy of $A$ itself to
$10.72$ percent, and lowers the accuracy of $B$ to $23.79$ percent. FGSM built on
$A$ gives $14.91$ on $A$ and $25.96$ on $B$. So the attack transfers, but
it weakens: about half of its strength is lost when it crosses models.

The most economical explanation is still the one of this chapter: at a small scale
the two models learn functions that are almost linear, and a direction that raises
the error on one linear function also raises it on another linear function whose
normal is close by. That is why transfer works best for one-step attacks. Our numbers
are consistent: the gap between white box and transfer is smaller for FGSM ($11$
points) than for PGD of 40 steps ($13$ points).

For a black-box attacker, transfer offers a cheap path: train a substitute model,
attack that model, then send the result to the victim model. Another route is to
estimate the gradient from differences of output over many queries, which is accurate
but expensive and easy to detect. The black box is weaker because its information is
less, not because its attacker is less clever.

The goal of the attacker also changes the numbers drastically. An untargeted attack
only has to make the label change; a targeted attack has to make the label become a
particular class. At $\varepsilon=0.15$, our PGD of 40 steps reaches the target on
$33.00$ percent of the examples, whereas the untargeted attack succeeds on
$89.06$ percent of the examples (accuracy is left at $10.72$ percent). With
$\varepsilon=0.30$, the target is reached on $99.50$ percent of the
examples. That is why the number ``the attack succeeds 90 percent'' without a note on
whether it is targeted cannot be compared with another number. For a defense
evaluator, a targeted test is harder for the attacker, whereas an untargeted test is
harder for the defender.

\section{Defenses: what works and what only appears to work}
\tingkat{advanced}

The defender has three families of answers. The first changes the data: add
adversarial examples to the training. The second changes the model: make its surface
smoother, for example with quantized inputs or smoothed outputs. The third adds tools
outside the model: detect strange inputs, or refuse answers that are not confident.
Only the first family has a clear basis, namely the saddle formulation above; the
second and third families often give good numbers for the wrong reason.

The most frequently repeated mistake has a name.

\begin{definisi}[Gradient obfuscation]
A defense performs \emph{gradient obfuscation} when it makes the gradient that the
attacker uses useless: missing, pointing the wrong way, or random. Such a defense
lowers the success of gradient-based attacks without raising the distance to the
decision boundary.
\end{definisi}

In everyday language: that guard does not become more observant, he starts using a
form that the attacker cannot read. Athalye, Carlini, and Wagner distinguish three
ways obfuscation happens. First, a \emph{shattered} gradient: an operation that is
nondifferentiable or numerically unstable makes the gradient absent or wrong. Second,
a \emph{random} gradient: a defense that adds randomness makes one sample of the
gradient not the true gradient. Third, a gradient that \emph{explodes or vanishes}: a
defense that evaluates the network many times in sequence, as in repeated
quantization, behaves like a very deep network when it is differentiated.
\index{obfuscated gradient}

Why is obfuscation not a defense? Because it changes nothing that a persistent
attacker sees. If the gradient is useless, the attacker uses a substitute gradient
from a smoother model; if randomness is the obstacle, he averages the gradient over
many samples; if the gradient explodes, he uses a finite-difference approximation.

We build the small example ourselves. Take model $A$ and put input quantization in
front of it: every pixel is rounded to a multiple of $1/8$. The derivative of rounding
is zero almost everywhere, so an attacker who differentiates through that path gets a
zero vector and does not move. The result looks convincing: clean accuracy
$97.82$ percent, and after ``PGD of 40 steps'' its accuracy is $96.98$ percent. Almost
no damage. But if the attacker uses the gradient of the model without quantization at
the point that has already been quantized, its accuracy falls to $13.23$ percent, only
a little higher than $10.72$ percent on the model without quantization. That defense
does not close the road; it only closes the most naive road. The shape of the picture
is in Figure~\ref{fig:kuantisasi-gradien}.

\begin{figure}[htbp]
\centering
\begin{tikzpicture}[scale=1.05,font=\scriptsize]
  \fill[layerteal!8] (0,0.5) -- (4,2.9) -- (4,4) -- (0,4) -- cycle;
  \fill[accent!6] (0,0) -- (4,0) -- (4,2.9) -- (0,0.5) -- cycle;
  \draw[very thick,ink] (0,0.5)--(4,2.9);
  \draw[->,layergrey] (-0.2,0)--(4.35,0) node[right]{$x_{1}$};
  \draw[->,layergrey] (0,-0.2)--(0,4.35) node[above]{$x_{2}$};
  \node[font=\scriptsize\bfseries] at (2,4.68) {original model};
  \node[layerteal] at (0.78,3.55) {class $+1$};
  \node[accent] at (3.05,0.8) {class $-1$};
  \node[layergrey,rotate=31] at (2.75,2.95) {boundary};
  \draw[-{Latex},very thick,accent] (1.2,2.6)--(1.92,1.4);
  \node[accent,above left=-2pt] at (1.56,2.05) {$\boldsymbol{\delta}$};
  \fill[ink] (1.2,2.6) circle (1.6pt);
  \node[ink,above left=-2pt] at (1.2,2.6) {$\mathbf{x}$};
  \draw[accent,thick,fill=accent!35] (1.92,1.4) ++(135:0.07) -- ++(-45:0.14)
    -- ++(-135:0.14) -- ++(45:0.14) -- cycle;
  \node[accent,below right=-2pt] at (1.92,1.4) {$\mathbf{x}+\boldsymbol{\delta}$};
\end{tikzpicture}
\hfill
\begin{tikzpicture}[scale=1.05,font=\scriptsize]
  \fill[layerteal!8] (0,0.5) -- (4,2.9) -- (4,4) -- (0,4) -- cycle;
  \fill[accent!6] (0,0) -- (4,0) -- (4,2.9) -- (0,0.5) -- cycle;
  \foreach \g in {0.5,1,1.5,2,2.5,3,3.5}
    {\draw[layergrey!30,very thin] (\g,0)--(\g,4);
     \draw[layergrey!30,very thin] (0,\g)--(4,\g);}
  \draw[very thick,ink] (0,0.5)--(4,2.9);
  \draw[->,layergrey] (-0.2,0)--(4.35,0) node[right]{$x_{1}$};
  \draw[->,layergrey] (0,-0.2)--(0,4.35) node[above]{$x_{2}$};
  \node[font=\scriptsize\bfseries] at (2,4.68) {with rounding to $1/8$};
  \fill[ink] (1.2,2.6) circle (1.6pt);
  \node[ink,above left=-2pt] at (1.2,2.6) {$\mathbf{x}$};
  \draw[layerplum,thick] (1.2,2.6) circle (0.13);
  \node[layerplum,anchor=west] at (1.4,2.78) {$\nabla$ path $=\mathbf{0}$};
  \draw[-{Latex},very thick,dashed,layerplum] (1.2,2.6)--(1.92,1.4);
  \node[layerplum,anchor=east] at (1.44,1.72) {$\nabla$ substitute};
  \draw[layerplum,thick,fill=layerplum!35] (1.92,1.4) ++(135:0.07)
    -- ++(-45:0.14) -- ++(-135:0.14) -- ++(45:0.14) -- cycle;
\end{tikzpicture}
\caption{Input quantization on the same model. Left: the true gradient gives a
direction $\boldsymbol{\delta}$ that takes $\mathbf{x}$ across to the other side of
the boundary, so that its label changes. Right: rounding every feature to a multiple
of $1/8$ makes the derivative of the path zero at almost all points, so the reported
gradient arrow stops; the decision boundary itself does not shift, and the gradient
of the model without rounding still takes it across. The square marks the point after
the perturbation.}
\label{fig:kuantisasi-gradien}
\end{figure}

Athalye and colleagues examined nine white-box defenses at ICLR 2018. Seven of them
depend on gradient obfuscation, and adapted attacks succeeded in breaking six
completely and one partially, under the threat model each of them claimed. This is not
a small error in one paper; it is a recurring pattern. That is why there is a list of
symptoms that can be used before trusting a robustness number. If a one-step attack is
stronger than an iterative attack, something is wrong: an iterative attack should
always be stronger in the white box. If a black-box attack is stronger than a white-box
one, something is wrong: the black box is a smaller set of information. If an attack
without a bound on the perturbation does not reach one hundred percent success,
something is wrong. If a random search finds an adversarial example while the gradient
attack does not, something is wrong. If enlarging the budget does not raise success,
something is wrong.

\begin{tantangan}
A service refuses to answer when the absolute value of its score margin is below 0.50,
and answers only when it is confident enough. In the linear example of the previous
part, the clean margin is 1.500 and $\lVert\mathbf{d}\rVert_1=6$. An attack at
$\varepsilon=0.25$ pushes the margin to zero, exactly at the threshold of reversal, so
it is refused, and the service looks safe. Determine the smallest budget that keeps
the attacker served, that is the one that switches the refusal off on the wrong side.
Is that rule a defense or obfuscation?

Hint 1: the refusal looks only at the size of the margin, not at its sign.

Hint 2: solve $1.500-6\varepsilon=-0.500$.

If you get stuck, go back to the section Worked example: a linear model you can compute
by hand.
\end{tantangan}

What does a correct evaluation look like? It is not one number, but one protocol. Fix
the threat model first: norm, $\varepsilon$, white box or black box, targeted or
not. Use the strongest attacker you can afford, and adapt it to the defense: if the
defense quantizes, attack through a substitute gradient; if it adds randomness,
average over several samples; if it filters the input, attack that filter as well.
Report the accuracy at several $\varepsilon$, not at one point, and check that the
curve falls monotonically. Report more than one random seed and include transfer
attacks from another model. Report the cost: the number of backward passes per
example, the time, and the hardware. If you want to talk about guarantees, use
certification, not testing; that is another path that is not discussed here.

\begin{untukanak}
Think about a multiple-choice exam that is graded by machine. The machine reads the
answer from the position of the bubble that is filled in. If we test the machine only
with answers that are filled in neatly, the machine will always be right. Then comes a
student who knows how the machine reads: he fills one other bubble a little more
thickly, and the machine reads it as a different answer. The paper is still the same
paper, the writing is still the same writing.

Now the interesting part. One school tries to fix that machine by scanning the paper
twice and taking the answer that appears most often. The same student just fills two
bubbles, and the machine can still be fooled. Another school adds a ruler so that the
position of the bubbles shifts a little; the patient student just tries a few
positions. All those additions make the cheating more troublesome, but they do not
make the machine able to tell an honest answer.

In this chapter we will find out why making things harder for the attacker is not the
same as closing the hole, and how to test a fix so that we are not fooled by a fix
that only appears to work.
\end{untukanak}

\section{Worked example: a linear model you can compute by hand}
\tingkat{intermediate}

Now we leave the networks and return to a model that can be computed without a
computer: a linear classifier with two classes and four features. Every number in this
part can be checked with a calculator.

Let class $A$ use the weights $\mathbf{w}_A=(1;-2;0.5;1.5)$ and class $B$
use $\mathbf{w}_B=(0;1;0.5;-0.5)$, without a bias. The example that is examined is
$\mathbf{x}=(1;0.5;1;1)$. The scores of the two classes are
\[
  s_A=\mathbf{w}_A^{\top}\mathbf{x}=1-1+0.5+1.5=2.0,
  \qquad
  s_B=\mathbf{w}_B^{\top}\mathbf{x}=0+0.5+0.5-0.5=0.5 .
\]
Its margin is $m=s_A-s_B=1.5$, so the model is confident but not excessively so. The
difference of the weights is $\mathbf{d}=\mathbf{w}_A-\mathbf{w}_B=(1;-3;0;2)$, with
$\lVert\mathbf{d}\rVert_1=1+3+0+2=6$ and
$\lVert\mathbf{d}\rVert_2=\sqrt{1+9+0+4}=\sqrt{14}\approx3.7417$.

An untargeted attack has to lower the margin. By the lemma, the largest decrease
is attained by $\boldsymbol{\delta}=-\varepsilon\,\mathrm{sign}(\mathbf{d})
=\varepsilon(-1;1;0;-1)$, and the size of the decrease is
$\varepsilon\lVert\mathbf{d}\rVert_1=6\varepsilon$. So the margin after the attack
is $1.5-6\varepsilon$, and the label changes when
\[
  \varepsilon>\varepsilon^{*}=\frac{1.5}{6}=0.25 .
\]
At $\varepsilon=0.25$ the margin is exactly zero, so the two classes tie and
$\arg\max$ chooses the first class; a clear reversal happens only a little above
$0.25$. The following table contains the whole course of it, and every row can be
checked by multiplying $\mathbf{d}$ by $\mathbf{x}+\boldsymbol{\delta}$.

\begin{center}
\begin{tabular}{rccc}
\hline
$\varepsilon$ & $s_A$ & $s_B$ & margin \\
\hline
$0$    & $2.000$ & $0.500$ & $+1.500$ \\
$0.05$ & $1.775$ & $0.575$ & $+1.200$ \\
$0.10$ & $1.550$ & $0.650$ & $+0.900$ \\
$0.15$ & $1.325$ & $0.725$ & $+0.600$ \\
$0.20$ & $1.100$ & $0.800$ & $+0.300$ \\
$0.25$ & $0.875$ & $0.875$ & $0.000$ \\
$0.30$ & $0.650$ & $0.950$ & $-0.300$ \\
\hline
\end{tabular}
\end{center}

If the attacker is bounded by the $\ell_2$ norm instead of $\ell_\infty$, the answer
changes. The maximum of $\mathbf{d}^{\top}\boldsymbol{\delta}$ on the $\ell_2$ ball of
radius $\varepsilon$ is $\varepsilon\lVert\mathbf{d}\rVert_2$, attained at
$\boldsymbol{\delta}=\varepsilon\mathbf{d}/\lVert\mathbf{d}\rVert_2$. So
$\varepsilon^{*}=1.5/\sqrt{14}\approx0.4009$. The $\ell_2$ budget that gives damage of
the same size as $\varepsilon=0.25$ on $\ell_\infty$ is
$0.4009$; both move the example to the same point, namely the point of intersection
$\mathbf{x}-\varepsilon\,\mathrm{sign}(\mathbf{d})$. This is a small example of a
general rule: the number $\varepsilon$ cannot be compared across norms.

Now the part that is often asked: is PGD stronger on this model? Every PGD step adds a
decrease of the margin of $\alpha\lVert\mathbf{d}\rVert_1
=6\alpha$, but the whole trajectory stays inside the ball, so the total decrease never
exceeds $6\varepsilon$. For a linear model, the worst point inside the ball really is
the point that FGSM chooses, so iteration cannot beat it. Hand computation shows it:
with $\varepsilon=0.10$ and 40 steps of
$\alpha=0.01$, the final margin $0.900$ is exactly the same as the result of FGSM at
$\varepsilon=0.10$.

\kertas{Repeat the example above for $\mathbf{w}_A=(2;1;-1;0.5)$,
$\mathbf{w}_B=(0;1;1;-0.5)$, and $\mathbf{x}=(1;1;0.5;1)$. Compute $s_A$, $s_B$,
$\mathbf{d}$, $\lVert\mathbf{d}\rVert_1$, and $\varepsilon^{*}$ for the
$\ell_\infty$ norm. Then compute $\varepsilon^{*}$ for the $\ell_2$ norm and explain
why the two numbers differ. Check with a three-line Python snippet. After that decide,
without running a program, whether PGD of 20 steps with
$\alpha=\varepsilon/10$ can change the label when $\varepsilon$ is only half of
$\varepsilon^{*}$.}

\section{Computational lab}
\tingkat{advanced}

This lab trains three small networks on handwritten digit data provided by the
\texttt{scikit-learn} library: 1797 examples, 64 features, values $0$ to $16$ which we
divide by $16$ so that they fall in $[0,1]$. Twelve hundred examples are used for
training and 597 for test. Model $A$ has one hidden layer of 64 units, model $B$ 128
units, and model $D$ is the same as $A$ but trained with FGSM at $\varepsilon=0.15$
(the reason for using $D$ rather than $A$ is that the comparison of clean and robust
accuracy should come from models of the same size).

All the gradients are computed by hand with NumPy, without an automatic differentiation
library, so that every step can be checked. The function \texttt{grad\_x} is the
backward pass of Chapter 25 that stops at the input, not at the weights.

\begin{lab}
\begin{verbatim}
import numpy as np, time
from sklearn.datasets import load_digits

X, y = load_digits(return_X_y=True)
X = X / 16.0                        # 0..16 -> 0..1
perm = np.random.default_rng(0).permutation(len(X))
tr, te = perm[:1200], perm[1200:]
Xtr, ytr, Xte, yte = X[tr], y[tr], X[te], y[te]
K = 10
Ytr, Yte = np.eye(K)[ytr], np.eye(K)[yte]

def init(dims, seed):
    r = np.random.default_rng(seed); W, b = [], []
    for i in range(len(dims) - 1):
        W.append(r.normal(0, np.sqrt(2.0 / dims[i]),
                          (dims[i], dims[i + 1])))
        b.append(np.zeros(dims[i + 1]))
    return W, b

def fwd(W, b, X):
    acts, pre, h = [X], [], X
    for i in range(len(W) - 1):
        z = h @ W[i] + b[i]; pre.append(z)
        h = np.maximum(z, 0); acts.append(h)
    z = h @ W[-1] + b[-1]; pre.append(z)
    z = z - z.max(1, keepdims=True); p = np.exp(z)
    return p / p.sum(1, keepdims=True), acts, pre

def loss_grad(W, b, X, Y):
    p, acts, pre = fwd(W, b, X); n = len(X)
    loss = -np.mean(np.sum(Y * np.log(p + 1e-12), 1))
    d = p - Y; gW = [None] * len(W); gb = [None] * len(W)
    for i in range(len(W) - 1, -1, -1):
        gW[i] = acts[i].T @ d / n + 1e-4 * W[i]; gb[i] = d.mean(0)
        if i > 0:
            d = (d @ W[i].T) * (pre[i - 1] > 0)
    return loss, gW, gb

def grad_x(W, b, X, Y):             # gradient of the error wrt the input
    p, acts, pre = fwd(W, b, X)
    d = (p - Y) / len(X)
    for i in range(len(W) - 1, 0, -1):
        d = (d @ W[i].T) * (pre[i - 1] > 0)
    return d @ W[0].T

def train(dims, seed, Xa, Ya, epochs=350, lr=0.01, adv_eps=None):
    W, b = init(dims, seed)
    mW = [np.zeros_like(w) for w in W]
    vW = [np.zeros_like(w) for w in W]
    mb = [np.zeros_like(v) for v in b]
    vb = [np.zeros_like(v) for v in b]
    for t in range(1, epochs + 1):
        if adv_eps is None:
            Xb = Xa
        else:
            g = grad_x(W, b, Xa, Ya)
            Xb = np.clip(Xa + adv_eps * np.sign(g), 0, 1)
        _, gW, gb = loss_grad(W, b, Xb, Ya)
        for i in range(len(W)):
            mW[i] = 0.9 * mW[i] + 0.1 * gW[i]
            vW[i] = 0.999 * vW[i] + 0.001 * gW[i] ** 2
            mb[i] = 0.9 * mb[i] + 0.1 * gb[i]
            vb[i] = 0.999 * vb[i] + 0.001 * gb[i] ** 2
            mh = mW[i] / (1 - 0.9 ** t)
            vh = vW[i] / (1 - 0.999 ** t)
            W[i] -= lr * mh / (np.sqrt(vh) + 1e-8)
            mh = mb[i] / (1 - 0.9 ** t)
            vh = vb[i] / (1 - 0.999 ** t)
            b[i] -= lr * mh / (np.sqrt(vh) + 1e-8)
    return W, b

def acc(W, b, X, y):
    return float(np.mean(fwd(W, b, X)[0].argmax(1) == y))

def fgsm(W, b, X, y, eps):
    g = grad_x(W, b, X, np.eye(K)[y])
    return np.clip(X + eps * np.sign(g), 0, 1)

def pgd(W, b, X, y, eps, steps, Ylab=None, arah=1):
    alpha = eps / 10
    r = np.random.default_rng(1)
    xa = np.clip(X + r.uniform(-eps, eps, X.shape), 0, 1)
    Y = np.eye(K)[y] if Ylab is None else Ylab
    for _ in range(steps):
        g = arah * np.sign(grad_x(W, b, xa, Y))
        xa = np.clip(np.clip(xa + alpha * g, X - eps, X + eps), 0, 1)
    return xa
\end{verbatim}
\end{lab}

The rest is calls and printing. We train $A$, $B$, and $D$, then sweep
$\varepsilon$ from $0.02$ to $0.30$.

\begin{verbatim}
WA, bA = train([64, 64, 10], 1, Xtr, Ytr)
WB, bB = train([64, 128, 10], 2, Xtr, Ytr)
WD, bD = train([64, 64, 10], 3, Xtr, Ytr, epochs=120, adv_eps=0.15)

print(f"clean test accuracy: A {acc(WA,bA,Xte,yte):.4f}"
      f"  B {acc(WB,bB,Xte,yte):.4f}  D {acc(WD,bD,Xte,yte):.4f}")
print(f"{'eps':>5} {'A clean':>9} {'A FGSM':>7} {'A PGD40':>8}"
      f" {'D FGSM':>7} {'D PGD40':>8}")
for eps in [0.02, 0.05, 0.10, 0.15, 0.20, 0.30]:
    xf = fgsm(WA, bA, Xte, yte, eps)
    xp = pgd(WA, bA, Xte, yte, eps, 40)
    xd = fgsm(WD, bD, Xte, yte, eps)
    xq = pgd(WD, bD, Xte, yte, eps, 40)
    print(f"{eps:5.2f} {acc(WA,bA,Xte,yte):9.4f} {acc(WA,bA,xf,yte):7.4f}"
          f" {acc(WA,bA,xp,yte):8.4f} {acc(WD,bD,xd,yte):7.4f}"
          f" {acc(WD,bD,xq,yte):8.4f}")

r = np.random.default_rng(7)
rs = np.sign(r.normal(size=Xte.shape))          # control: random signs
print("random signs eps=0.15:",
      f"{acc(WA,bA,np.clip(Xte+0.15*rs,0,1),yte):.4f}")
print("PGD steps, eps=0.15, alpha=eps/10, random start")
for k in [1, 2, 5, 10, 20, 40]:
    print(f"  {k:2d} steps:",
          f"{acc(WA,bA,pgd(WA,bA,Xte,yte,0.15,k),yte):.4f}")
print("  FGSM (1 step, alpha=eps):",
      f"{acc(WA,bA,fgsm(WA,bA,Xte,yte,0.15),yte):.4f}")
\end{verbatim}

The actual output, as it stands, for the block above:

\begin{verbatim}
clean test accuracy: A 0.9799  B 0.9849  D 0.9648
  eps  A clean  A FGSM  A PGD40  D FGSM  D PGD40
 0.02    0.9799  0.9548   0.9548  0.9497   0.9497
 0.05    0.9799  0.8476   0.8425  0.9162   0.9129
 0.10    0.9799  0.4992   0.4590  0.8241   0.8174
 0.15    0.9799  0.1491   0.1072  0.6667   0.6265
 0.20    0.9799  0.0184   0.0050  0.4858   0.3719
 0.30    0.9799  0.0000   0.0000  0.0687   0.0201
random signs eps=0.15: 0.9548
PGD steps, eps=0.15, alpha=eps/10, random start
   1 steps: 0.9464
   2 steps: 0.8878
   5 steps: 0.6667
  10 steps: 0.2479
  20 steps: 0.1122
  40 steps: 0.1072
  FGSM (1 step, alpha=eps): 0.1491
\end{verbatim}

The first reading of that table: clean accuracy $97.99$ percent, and at
$\varepsilon=0.15$ only $14.91$ percent is left.

The second reading: the gap between FGSM and PGD widens in the middle, from zero at
$0.02$ to $4.19$ points at $0.15$. At a small budget one step already touches the
wall of the box; at a large budget several steps are needed to find a worse
direction.

The third reading, and the most important one for this chapter: the row of random
signs. If a perturbation of size $\varepsilon=0.15$ is given with random signs on
every pixel, accuracy only falls to $95.48$ percent. FGSM with the same size lowers it
to $14.91$ percent. So the choice of signs by the gradient is worth about eighty
points of accuracy: it is the direction that works, not the size.

The fourth reading: at $\varepsilon=0.30$, $A$ is already at zero percent while $D$ is
still at $6.87$ percent under FGSM and $2.01$ percent under PGD. Adversarial training
does not make the model immune; it shifts the curve to the right.

The next block measures two things that are not visible in the table: whether the
gradient we compute really is the gradient, and how much it costs.

\begin{verbatim}
i0 = int(np.argmax(np.abs(grad_x(WA, bA, Xte, Yte)).sum(1)))
x0, y0 = Xte[i0:i0 + 1], Yte[i0:i0 + 1]
g0 = grad_x(WA, bA, x0, y0)
L1 = lambda z: loss_grad(WA, bA, z, y0)[0]
print(f"finite difference at sample i={i0}, h=1e-6")
for j in np.argsort(-np.abs(g0[0]))[:3]:
    e = np.zeros_like(x0); e[0, j] = 1e-6
    fd = (L1(x0 + e) - L1(x0 - e)) / 2e-6
    print(f"  j={j:2d} analytic {g0[0,j]:+.6e} finite-diff {fd:+.6e}"
          f" error {abs(fd-g0[0,j])/abs(g0[0,j]):.0e}")
g_all = grad_x(WA, bA, Xte, Yte)
lin = 0.15 * np.abs(g_all).sum(1).mean()
lc = loss_grad(WA, bA, Xte, Yte)[0]
la = loss_grad(WA, bA, fgsm(WA, bA, Xte, yte, 0.15), Yte)[0]
print(f"mean |g|_1 {np.abs(g_all).sum(1).mean():.4f};"
      f" linear prediction {lin:.4f} nats; actual {la-lc:.3f} nats")

t0 = time.perf_counter(); fgsm(WA, bA, Xte, yte, 0.15)
t1 = time.perf_counter(); pgd(WA, bA, Xte, yte, 0.15, 40)
t2 = time.perf_counter()
print(f"time for 597 examples: FGSM {t1-t0:.4f} s  PGD40 {t2-t1:.4f} s")
\end{verbatim}

Its output:

\begin{verbatim}
finite difference at sample i=145, h=1e-6
  j=53 analytic -6.391833e+00 finite-diff -6.391833e+00 error 2e-11
  j=38 analytic +6.211518e+00 finite-diff +6.211518e+00 error 1e-10
  j=13 analytic -6.175321e+00 finite-diff -6.175321e+00 error 8e-11
mean |g|_1 0.0056; linear prediction 0.0008 nats; actual 4.862 nats
time for 597 examples: FGSM 0.0006 s  PGD40 0.0302 s
\end{verbatim}

The first three lines confirm that our backward pass is correct: the analytic gradient
and the finite difference with $h=10^{-6}$ agree to about ten digits, with relative
errors $10^{-10}$ and $10^{-11}$. For a network that uses ReLU, agreement that smooth
is reasonable: at the scale $10^{-6}$ no unit changes sign, so the function really is
linear around that point. If these numbers differed a lot, what would be wrong is the
backward-pass code, not the attack.

The fourth line is the surprising part. The average $\lVert\nabla L\rVert_1$ is only
$0.0056$, so the linear prediction for $\varepsilon=0.15$ is $0.0008$
nats. The actual rise in the error is $4.862$ nats. So the linear term
underestimates the damage by about four thousand times. This is not a sign that the
gradient is wrong (the three lines above prove the opposite), but a sign that $L$ is
strongly curved at the distance $0.15$: ReLU units change sign, and the first-order
approximation loses its meaning. What remains useful is the sign, as the row of random
signs showed.

The last line is about time: for 597 examples, FGSM needs $0.6$ milliseconds and
PGD of 40 steps $30$ milliseconds, about fifty times more expensive. This number
depends on the machine and on vectorization, so what is worth remembering is its
order of magnitude, not the exact number.

Finally, a demonstration of gradient obfuscation. We add input quantization to
multiples of $1/8$ in front of model $A$, then attack in two ways.

\begin{verbatim}
def q(z): return np.round(z * 8) / 8
accq = lambda W, b, Z, yy: float(np.mean(fwd(W, b, q(Z))[0]
                                        .argmax(1) == yy))
def pgd_naif(W, b, Z, yy, eps, steps):
    r = np.random.default_rng(1)
    z = np.clip(Z + r.uniform(-eps, eps, Z.shape), 0, 1)
    for _ in range(steps):          # gradient = zero, so it stands still
        z = np.clip(np.clip(z, Z - eps, Z + eps), 0, 1)
    return z
def pgd_bpda(W, b, Z, yy, eps, steps):
    r = np.random.default_rng(1)
    z = np.clip(Z + r.uniform(-eps, eps, Z.shape), 0, 1)
    for _ in range(steps):
        g = np.sign(grad_x(W, b, q(z), np.eye(K)[yy]))
        z = np.clip(np.clip(z + eps / 10 * g, Z - eps, Z + eps), 0, 1)
    return z
print("inputs quantized to multiples of 1/8")
print(f"  clean {accq(WA,bA,Xte,yte):.4f}"
      f"  PGD40 naive {accq(WA,bA,pgd_naif(WA,bA,Xte,yte,0.15,40),yte):.4f}"
      f"  PGD40 BPDA"
      f" {accq(WA,bA,pgd_bpda(WA,bA,Xte,yte,0.15,40),yte):.4f}")
\end{verbatim}

Its output:

\begin{verbatim}
inputs quantized to multiples of 1/8
  clean 0.9782  PGD40 naive 0.9698  PGD40 BPDA 0.1323
\end{verbatim}

This is the most concise form of a false defense. Its clean number is only $0.17$
points below the original model, and the naive attack does not move at all: if that
were all that was reported, this defense would look almost perfect. The substitute
gradient brings it down to $13.23$ percent, almost the same as $10.72$ percent on the
model without quantization. What changes above all is the readability of its gradient.

\section{Cost, limitations, and how to report}
\tingkat{basic}

The computational budget determines which attack is honestly used. FGSM requires one
backward pass per example. PGD of $k$ steps requires $k$ passes. In our experiment, the
cost of PGD of 40 steps is about fifty times that of FGSM, for a result that, at small
$\varepsilon$, is almost the same. For the assessment of a defense, that cost is
reasonable: a robustness number measured with a weak attack is meaningless. For
training, that cost multiplies again: every epoch demands a new attack.

Adversarial training also has an accuracy price. Model $D$ loses $1.51$ points of
clean accuracy in exchange for a rise of about fifty points at $\varepsilon=0.15$.
Madry and colleagues report that larger models give better robustness, which shows
that robustness is not a free property of the architecture, but capacity that has to
be paid for.

There are three limitations before the numbers in this chapter are used.

First, scale. Our data has only 64 features and 10 classes, with 597 test examples. On
large images, a meaningful budget is much smaller than $0.15$; on ImageNet,
Goodfellow and colleagues used $0.007$ and already flipped the classification. The
number $\varepsilon$ cannot be moved between datasets, because its meaning depends on
the scale of the features. On our data, the pixels are already quantized to multiples
of $1/16$, so that $\varepsilon=0.05$ means less than one quantization step per
pixel.

Second, the coverage of attacks. What we run is only two families: FGSM and PGD with
$\ell_\infty$. A model that resists both is not necessarily resistant to an
$\ell_2$ attack, to an attack that changes a single pixel, or to an attack that sticks
an object into the real world. Every robustness claim holds only for the family of
attacks that was tested, and that is why the evaluation protocol in the previous part
demands several kinds of attack at once.

Third, there is no guarantee. All the numbers here are results of testing, not a
proved lower bound. Model $D$, which reaches $62.65$ percent under PGD of 40 steps,
may still be broken by a better attack; we do not know, because what was not found is
not proof that it does not exist. For a stronger statement, certification is needed:
for example with randomized smoothing or with formal verification on a small network.
Both give a guarantee on a certain radius, at the price of far higher accuracy and
computation cost, and both are outside the scope of this chapter.

An honest way of reporting contains five things: the complete threat model (norm,
$\varepsilon$, the knowledge of the attacker, targeted or not), the attacker used
together with the number and size of its steps, the clean accuracy and the accuracy at
several $\varepsilon$, the number of seeds and the spread of the results, and the
computational cost. Without that, the number ``90 percent robust'' cannot be checked
and cannot be compared. Chapter 27 has already taught us that metrics can deceive;
here the metric deceives in a subtler way, because its denominator is right and only
its numerator is chosen.

\begin{untukdosen}
The teaching order that has worked: start from the table of accuracy against
$\varepsilon$, not from the definition. Students see $97.99$ and $14.91$ on the same
model, then ask why. Only after that come the definition of the threat model, the
lemma on the maximum of a linear form, and FGSM. Only after FGSM runs do we enter PGD
and the saddle formulation. The reverse order, that is definitions first and then
attacks, makes students memorize notation without ever seeing the damage.

Two typical misconceptions. First, students think an adversarial attack is the same as
adding noise, so they conclude that a model resistant to noise will be resistant to
attacks. The row of random signs in the lab refutes that with one number: noise of
size $\varepsilon=0.15$ only lowers accuracy to $95.48$ percent, whereas FGSM with the
same size lowers it to $14.91$ percent. Second, they think one accuracy number under
an attack is already evidence of robustness. The quantization demonstration in the lab
shows the opposite: $96.98$ percent looks safe until someone uses a substitute
gradient and gets $13.23$ percent.

A cheap blackboard activity: write the four-feature linear example from the worked
example part, then ask the class to compute $\varepsilon^{*}$ for $\ell_\infty$ and
$\ell_2$. After that ask them to predict, without computing, whether adding PGD steps
changes the result on that linear model. The correct answer, together with its reason,
forces them to understand that iteration is useful only because $f$ is not linear.

A sign of readiness to go on: students can explain why only the sign of the gradient
is used, can compute $\varepsilon^{*}$ for a linear model, can name one symptom of
gradient obfuscation, and can design one attack adapted to a defense that quantizes
the input. If they demand numbers at several $\varepsilon$ before believing, they are
ready to enter Chapter 30.
\end{untukdosen}

\section{Summary}

\begin{ringkasan}
\begin{itemize}
  \item A robustness claim is always relative to the threat model: the norm, the
        budget $\varepsilon$, the knowledge of the attacker, and whether the goal is
        targeted or not. Without that, two numbers cannot be compared.
  \item FGSM shifts every coordinate by $\varepsilon$ following the sign of the
        gradient of the error. On a linear model it is exact and optimal, because
        $\max_{\lVert\boldsymbol{\delta}\rVert_\infty\le\varepsilon}
        \mathbf{w}^{\top}\boldsymbol{\delta}=\varepsilon\lVert\mathbf{w}\rVert_1$.
  \item Linearization justifies the direction, not the distance. In our experiment the
        linear prediction is $0.0008$ nats, whereas the actual rise is $4.862$ nats;
        what remains useful is the pattern of signs.
  \item PGD adds iteration and projection onto the box. At $\varepsilon=0.15$
        accuracy falls from $94.64$ percent (1 small step) to $10.72$ percent
        (40 steps), and is flat after about twenty steps.
  \item Adversarial training is an approximation of the saddle problem
        $\min_\theta\mathbb{E}\max_\delta L$. It raises robustness at the price of
        clean accuracy: our model loses $1.51$ points to rise from $10.72$
        to $62.65$ percent.
  \item Gradient obfuscation lowers the success of attacks without raising the distance
        to the decision boundary. Our example: $96.98$ percent under the naive
        attack, $13.23$ percent under the attack with a substitute gradient.
  \item An honest evaluation uses the strongest attack that can be afforded, adapted
        to the defense, at several $\varepsilon$ and several seeds, with the cost
        reported.
\end{itemize}
\end{ringkasan}

\section{Exercises}

\begin{latihan}
\begin{enumerate}
  \item Compute $\varepsilon^{*}$ for the linear model with
        $\mathbf{w}_A=(1;-1;2)$, $\mathbf{w}_B=(0;1;1)$ and
        $\mathbf{x}=(1;2;1)$ under the $\ell_\infty$ norm, then under the $\ell_2$ norm.
        Check both answers with Python. (Hint: compute
        $\mathbf{d}=\mathbf{w}_A-\mathbf{w}_B$ first, then compare
        $\lVert\mathbf{d}\rVert_1$ with $\lVert\mathbf{d}\rVert_2$.)

  \item Prove that $\max_{\lVert\boldsymbol{\delta}\rVert_2\le\varepsilon}
        \mathbf{w}^{\top}\boldsymbol{\delta}=\varepsilon\lVert\mathbf{w}\rVert_2$,
        and determine all $\boldsymbol{\delta}$ that attain the maximum. Explain
        why the answer has a different form from the $\ell_\infty$ case.

  \item Prove that on a linear classifier, PGD with any number of
        steps never produces an error larger than FGSM
        with the same $\varepsilon$. (Hint: the error is a linear function of
        $\boldsymbol{\delta}$ on a convex set; the maximum of a linear function on a
        compact convex set is attained at an extreme point. Show that an extreme point
        of the $\ell_\infty$ ball has the form $\pm\varepsilon$ in every coordinate,
        then compare with the FGSM point.)

  \item Assess this claim: \emph{because a model reaches accuracy $98$ percent on
        test data and stays at $95$ percent under random perturbation of size
        $\varepsilon$, that model is resistant to adversarial attacks.} Decide true,
        false, or not enough evidence, and name one number from this chapter that
        supports your decision.

  \item From the lab table, compute the success rate of the untargeted attack for
        model $A$ at $\varepsilon=0.05$, $0.10$, and $0.15$, then for
        model $D$ at the same budgets. State the results in percent and
        explain why the gap between $A$ and $D$ narrows at large $\varepsilon$.

  \item Show that the projection onto
        $\{\mathbf{y}:\lVert\mathbf{y}-\mathbf{x}\rVert_\infty\le\varepsilon\}
        \cap[0,1]^{d}$ may be done by clipping first to the $\varepsilon$ box and then
        to $[0,1]$, or the other way round, and the result is the same.
        (Hint: both sets are products of intervals, so it is enough to check
        one coordinate.)

  \item On model $D$, the clean accuracy is $96.48$ percent and the accuracy under PGD
        of 40 steps at $\varepsilon=0.15$ is $62.65$ percent. On model $A$,
        the numbers are $97.99$ and $10.72$. Compute the price of robustness in
        points of clean accuracy per point of robustness gained, then explain why that
        ratio is not a complete measure.

  \item A defense rounds the probability output to $0$ or $1$
        before the decision is taken, so that the gradient with respect to the input
        becomes zero. Design one evaluation protocol to test its robustness claim, and
        name which symptom from the list in this chapter is most easily seen if that
        defense only obfuscates the gradient.

  \item Given two models with clean accuracy $97.99$ and $98.49$ percent.
        At $\varepsilon=0.15$, PGD of 40 steps built on the first model lowers
        the accuracy of the first model to $10.72$ percent and of the second model to
        $23.79$ percent. Compute the gap between the success of the white-box attack
        and of the transfer attack, then explain why that gap is not evidence that the
        second model is safer. (Hint: an attack built specifically for the second model
        is also available; see row B in the lab output.)

  \item Explain why a targeted attack is harder than an untargeted
        attack, then use the numbers $33.00$ percent at $\varepsilon=0.15$ and
        $99.50$ percent at $\varepsilon=0.30$ to show that the difficulty
        depends on the budget, not only on the kind of attack.
\end{enumerate}
\end{latihan}

\begin{untukpeneliti}
What is proved in this chapter is little. The lemma on the maximum of a linear form is
a one-line proof. The linearization inequality with remainder
$M\lVert\boldsymbol{\delta}\rVert_2^2/2$ is the mean value theorem in integral form,
valid for every differentiable function with a Lipschitz gradient. The exactness of
FGSM on logistic regression is a result of Goodfellow, Shlens, and Szegedy (ICLR 2015),
the section ``Adversarial Training of Linear Models versus Weight Decay''. The saddle
formulation and the claim that PGD can be relied on come from Madry, Makelov, Schmidt,
Tsipras, and Vladu (ICLR 2018), sections 2 and 3; the latter is an empirical claim,
not a theorem, and its own authors call the inner problem nonconcave.

What is more often popularized than proved is the claim that linearity explains
adversarial examples. As an argument about direction it is useful; as an explanation of
the size of the damage it fails in our experiment: a prediction of $0.0008$ nats
against an actual $4.862$ nats, with a gradient that has been verified to $10^{-10}$.
What works is the pattern of signs, and that is why FGSM beats random signs
($14.91$ against $95.48$ percent). A second claim that is exaggerated: adversarial
training gives robustness to the attacker and the budget used during training, at the
price of clean accuracy, without a guarantee for other attacks.

Limits of validity: 64 features, 597 test examples, one hidden layer, one seed, two
families of $\ell_\infty$ attacks, without certification. The numbers here demonstrate
a mechanism, not a ranking of robustness. Open questions: how large a radius can be
certified without erasing the benefit in accuracy; why transfer works that well;
and what threat model makes sense for real systems such as intrusion detectors, which
is the subject of Chapter 30.
\end{untukpeneliti}

\peta{This chapter follows Goodfellow, Shlens, and Szegedy, \emph{Explaining and
Harnessing Adversarial Examples} (ICLR 2015), sections 3 (``The Linear Explanation of
Adversarial Examples''), 4 (``The Fast Gradient Sign Method''), and 5 (``Adversarial
Training of Linear Models versus Weight Decay'') for the derivation of FGSM, its
exactness on logistic regression, and the MNIST and CIFAR-10 numbers. The saddle
formulation, the PGD algorithm, and the MNIST and CIFAR-10 robustness numbers come from
Madry, Makelov, Schmidt, Tsipras, and Vladu, \emph{Towards Deep Learning Models
Resistant to Adversarial Attacks} (ICLR 2018), sections 2, 3, and 4, including
Algorithm 1. The discussion of gradient obfuscation follows Athalye, Carlini, and
Wagner, \emph{Obfuscated Gradients Give a False Sense of Security} (ICML 2018),
sections 2 to 4, including the three kinds of obfuscation and the five symptoms used
here. The phenomenon of transferability was first reported by Szegedy and colleagues,
\emph{Intriguing Properties of Neural Networks} (ICLR 2014). All the numbers in the
tables, the lab, and the worked example are our own computation with Python 3.14,
NumPy 2.5, and the \texttt{load\_digits} dataset from scikit-learn 1.9: Adam training
for 350 epochs, a train-test split of 1200/597 with seed 0, and input gradients
verified with finite differences.}
